% !TeX root = proyecto.tex

%========================================================================================
% REQUERIMIENTOS DEL SISTEMA — TT-2 / 2026-A135
% Distribución: 30% Casos de Uso (UI) · 70% Requerimientos del Motor Back-end
%========================================================================================

\chapter{Especificación de Requerimientos del Sistema}
\label{cap:requerimientos}

El presente capítulo formaliza el contrato técnico del sistema bajo el protocolo
2026-A135. Atendiendo a la naturaleza predominantemente autónoma del sistema, la
especificación refleja una distribución de peso de \textbf{30\% en Casos de Uso}
---interacciones del usuario con el dashboard de resultados--- y \textbf{70\% en
Requerimientos de Sistema}, que definen el comportamiento del motor de recolección,
procesamiento y análisis semántico que opera de forma completamente independiente
de la intervención humana. El usuario final consume únicamente los resultados que
el pipeline de IA procesó de manera autónoma; no interviene en ninguna etapa del
análisis.

%========================================================================================
\section{Casos de Uso del Sistema (CU)}
\label{sec:cu}
%========================================================================================

Los Casos de Uso modelan las interacciones entre los actores externos (investigadores,
analistas de OSC, coordinadores del DIF) y el dashboard web. La Tabla~\ref{tab:cu_resumen}
presenta el catálogo completo; los actores interactúan exclusivamente con resultados
pre-procesados por el pipeline de IA.

\begin{table}[H]
\centering
\caption{Catálogo de Casos de Uso del Sistema}
\label{tab:cu_resumen}
\renewcommand{\arraystretch}{1.4}
\begin{tabular}{|p{1.5cm}|p{3.8cm}|p{8.3cm}|}
\hline
\rowcolor{black}
\textcolor{white}{\textbf{ID}} &
\textcolor{white}{\textbf{Nombre}} &
\textcolor{white}{\textbf{Descripción}} \\
\hline
\textbf{CU-01} &
Autenticación de Usuario &
El actor introduce credenciales (usuario + contraseña). El sistema verifica el
hash Argon2id almacenado en PostgreSQL. En caso de éxito genera un token de
sesión criptográfico (Flask-Login) con TTL de 8 horas. Tras 5 intentos fallidos
se bloquea la cuenta. \\
\hline
\textbf{CU-02} &
Visualizar Panel de Estadísticas &
El actor accede al dashboard principal. El sistema recupera de PostgreSQL las métricas
de las últimas 24/48/72 horas: total de noticias procesadas, distribución por nivel
de prioridad (Alta/Media/Baja), cantidad de casos en seguimiento y tópicos BERTopic
activos. Todos los datos son resultado del último ciclo autónomo del pipeline. \\
\hline
\textbf{CU-03} &
Filtrar Noticias por Ejes Semánticos &
El actor selecciona filtros combinados: eje semántico (Feminicidio / NNA / Orfandad
/ Violencia), nivel de prioridad, entidad federativa o rango de fechas. El sistema
ejecuta una consulta FTS sobre el índice GIN (\texttt{to\_tsquery}) y retorna los
resultados ordenados por \texttt{ts\_rank} en \textless~200~ms. \\
\hline
\textbf{CU-04} &
Consultar Ficha de Caso Identificado &
El actor selecciona una noticia de prioridad Alta. El sistema renderiza la ficha de
caso mostrando: título, medio, fecha, resumen semántico, clúster BERTopic asignado,
score BETO, y las detecciones de víctimas y entidades geográficas identificadas por
el pipeline. El actor puede añadir notas de investigación. \\
\hline
\textbf{CU-05} &
Gestionar Fuentes de Datos (CRUD) &
El actor agrega, edita o desactiva una URL de fuente RSS o HTML. Al agregar, el
sistema ejecuta un health check HTTP GET y valida que el endpoint responda con
código 200. Las fuentes desactivadas son excluidas del siguiente ciclo de recolección
sin eliminar su historial. \\
\hline
\textbf{CU-06} &
Marcar Caso en Seguimiento &
El actor activa la bandera \texttt{en\_seguimiento} sobre un caso identificado.
El sistema persiste el registro en la tabla \texttt{detecciones} y lo añade a la
vista ``Seguimiento de Casos'' donde es monitoreable de forma independiente. \\
\hline
\textbf{CU-07} &
Exportar Datos Estructurados (CSV) &
El actor solicita la exportación de un subconjunto filtrado de noticias. El sistema
serializa los campos \texttt{titulo, medio, fecha, prioridad, cluster\_id, score\_beto,
victimas, entidades} en formato CSV y lo entrega como descarga HTTP con tipo MIME
\texttt{text/csv; charset=utf-8}. \\
\hline
\textbf{CU-08} &
Ejecutar Análisis Manual On-Demand &
El actor solicita un ciclo de análisis inmediato (sin esperar el scheduler de 6 horas).
El sistema encola la tarea de recolección+análisis y retorna un identificador de trabajo.
El actor puede consultar el estado del proceso mediante polling al endpoint
\texttt{/api/v1/status/\{job\_id\}}. \\
\hline
\end{tabular}
\end{table}

%========================================================================================
\section{Requerimientos Funcionales del Motor Back-end}
\label{sec:rf}
%========================================================================================

Los Requerimientos Funcionales (RF) de esta sección describen el comportamiento
interno del sistema que opera de forma autónoma, sin intervención del usuario.
Se organizan en cinco subsistemas del pipeline: Recolección, Deduplicación,
Análisis Semántico, Agrupamiento y Persistencia.

%----------------------------------------------------------------------------------------
\subsection{RF — Subsistema de Recolección de Datos}
\label{subsec:rf_recoleccion}
%----------------------------------------------------------------------------------------

\begin{table}[H]
\centering
\caption{Requerimientos Funcionales: Subsistema de Recolección}
\label{tab:rf_recoleccion}
\renewcommand{\arraystretch}{1.4}
\begin{tabular}{|p{1.5cm}|p{3.5cm}|p{8.5cm}|}
\hline
\rowcolor{black}
\textcolor{white}{\textbf{ID}} &
\textcolor{white}{\textbf{Requerimiento}} &
\textcolor{white}{\textbf{Descripción Técnica}} \\
\hline
\textbf{RF-01} &
Ingesta RSS Multi-fuente &
El sistema debe parsear feeds RSS 2.0 y Atom de 45~fuentes configuradas mediante
\texttt{feedparser}, extrayendo por cada \texttt{<item>}: \texttt{title},
\texttt{link}, \texttt{summary}, \texttt{published}. El intervalo de recolección
es de 6 horas, gestionado por \texttt{scheduler.py}. Los feeds no modificados
(mismo \texttt{ETag} o \texttt{Last-Modified}) son omitidos sin re-descarga. \\
\hline
\textbf{RF-02} &
Recolección HTML con Evasión de WAF &
Para fuentes sin RSS, el sistema debe realizar requests HTTP mediante
\texttt{StealthSession} v7.0, que utiliza \texttt{cloudscraper} para generar
hashes TLS/JA3 idénticos a Chrome/Windows. Los intervalos entre requests deben
seguir la distribución $\mathcal{LN}(\mu=2.5, \sigma=0.7)$ (mediana $\approx$12~s).
El sistema debe respetar las directivas \texttt{robots.txt} de cada dominio. \\
\hline
\textbf{RF-03} &
Circuit Breaker por Dominio &
Ante 3 fallos HTTP consecutivos sobre un mismo dominio (códigos 403, 429, 5xx),
el sistema debe abrir el Circuit Breaker y suspender requests a ese dominio por
300~s. El estado (abierto/cerrado/semi-abierto) debe persistirse en memoria de
proceso para sobrevivir fallos de red transitorios. \\
\hline
\textbf{RF-04} &
Recolección Histórica (Google News API) &
El sistema debe ejecutar consultas a Google News mediante 14 queries de dominio
predefinidas (ej.\ \textit{``feminicidio hijos huérfanos México''}) para recuperar
noticias históricas con antigüedad de hasta 90 días. Los resultados deben integrarse
al mismo pipeline de deduplicación y análisis que las noticias RSS. \\
\hline
\end{tabular}
\end{table}

%----------------------------------------------------------------------------------------
\subsection{RF — Subsistema de Deduplicación en Cascada}
\label{subsec:rf_dedup}
%----------------------------------------------------------------------------------------

\begin{table}[H]
\centering
\caption{Requerimientos Funcionales: Subsistema de Deduplicación}
\label{tab:rf_dedup}
\renewcommand{\arraystretch}{1.4}
\begin{tabular}{|p{1.5cm}|p{3.5cm}|p{8.5cm}|}
\hline
\rowcolor{black}
\textcolor{white}{\textbf{ID}} &
\textcolor{white}{\textbf{Requerimiento}} &
\textcolor{white}{\textbf{Descripción Técnica}} \\
\hline
\textbf{RF-05} &
Hash Exacto MD5 &
Antes de insertar cualquier noticia, el sistema debe calcular el hash MD5 del
contenido textual normalizado (minúsculas, sin stopwords) y verificar su ausencia
en la tabla \texttt{noticias.hash\_md5}. Las colisiones exactas se descartan en
$O(1)$ sin procesamiento adicional. \\
\hline
\textbf{RF-06} &
Detección de Cuasi-duplicados (SimHash) &
Para noticias que superan la barrera MD5, el sistema debe calcular un SimHash de
64~bits basado en pesos TF-IDF de los tokens. Dos noticias con distancia de Hamming
$\leq 6$ bits se consideran cuasi-duplicadas; la más reciente se descarta y la
más antigua se conserva. \\
\hline
\textbf{RF-07} &
Similitud de Jaccard sobre Títulos &
El sistema debe tokenizar el título de cada noticia nueva y calcular el coeficiente
de Jaccard contra los títulos del corpus reciente (ventana de 48~h). Un valor
$J \geq 0.70$ clasifica el par como redundante; la nueva noticia se descarta. \\
\hline
\textbf{RF-08} &
Similitud Coseno TF-IDF sobre Contenido &
Como última barrera, el sistema debe vectorizar el contenido de noticias candidatas
con TF-IDF y calcular la similitud coseno. Un valor $\cos \geq 0.85$ confirma la
duplicación semántica; solo se retiene el documento con mayor score BETO. \\
\hline
\end{tabular}
\end{table}

%----------------------------------------------------------------------------------------
\subsection{RF — Subsistema de Análisis Semántico Neuro-Simbólico}
\label{subsec:rf_nlp}
%----------------------------------------------------------------------------------------

\begin{table}[H]
\centering
\caption{Requerimientos Funcionales: Subsistema de Análisis Semántico (Motor BETO)}
\label{tab:rf_nlp}
\renewcommand{\arraystretch}{1.4}
\begin{tabular}{|p{1.5cm}|p{3.5cm}|p{8.5cm}|}
\hline
\rowcolor{black}
\textcolor{white}{\textbf{ID}} &
\textcolor{white}{\textbf{Requerimiento}} &
\textcolor{white}{\textbf{Descripción Técnica}} \\
\hline
\textbf{RF-09} &
Preprocesamiento de Texto &
El sistema debe normalizar cada noticia: eliminación de HTML residual, normalización
Unicode NFC, colapso de espacios y reducción a 512 tokens mediante truncación
centrada (primeros 256 + últimos 256 tokens WordPiece) para respetar el límite
de secuencia de BETO. \\
\hline
\textbf{RF-10} &
Bypass de Oro (Capa Simbólica Afirmativa) &
Antes de invocar BETO, el sistema debe evaluar el texto contra el diccionario de
``Palabras de Oro'' (expresiones de alta precisión: \textit{``hijos quedaron
huérfanos''}, \textit{``menores al cuidado del DIF''}, etc.). Si se activa al
menos un patrón, la noticia recibe prioridad Alta directamente sin inferencia
neuronal, reduciendo la carga computacional. \\
\hline
\textbf{RF-11} &
Bozal de Penalización (Capa Simbólica Restrictiva) &
El sistema debe evaluar patrones de penalización antes de la inferencia BETO
(e.g., \textit{``menor detenido''}, \textit{``estadística de feminicidios''},
\textit{``feminicidio de una menor''} [víctima directa menor de edad]). Cada
patrón activado aplica una penalización $\delta \in [0.1, 0.4]$ al score final.
La suma de penalizaciones no puede exceder 0.8 (límite de descarte). \\
\hline
\textbf{RF-12} &
Clasificación Zero-Shot BETO &
El sistema debe calcular el embedding BETO (\texttt{[CLS]} token, última capa)
del texto preprocesado y su similitud coseno contra los embeddings precomputados
de las descripciones de las clases ``relevante'' y ``no relevante''. La similitud
debe amplificarse con factor $\alpha=5$ (temperature scaling) antes de la función
softmax para polarizar distribuciones sub-confiadas. \\
\hline
\textbf{RF-13} &
Evaluación de 4 Ejes Semánticos &
La clasificación final debe desagregarse en 4 ejes independientes: (1)~Feminicidio
confirmado, (2)~Presencia de NNA en la narrativa, (3)~Indicador de orfandad/resguardo
institucional, (4)~Violencia de género contextual. Cada eje genera un score
$\in [0, 1]$ que contribuye al score compuesto mediante el Factor Alpha Dinámico:
$S_{final} = \alpha_{ctx} \cdot S_{BETO} + (1-\alpha_{ctx}) \cdot S_{h}$. \\
\hline
\textbf{RF-14} &
Asignación de Prioridad &
El sistema debe asignar nivel de prioridad según umbrales calibrados:
\textbf{Alta}: $S_{final} \geq 0.75$ o Bypass de Oro activado;
\textbf{Media}: $0.50 \leq S_{final} < 0.75$;
\textbf{Baja}: $S_{final} < 0.50$ (registrada pero no alertada).
Las noticias con Bozal acumulado $\geq 0.8$ se descartan y se registran
en el log de descarte para auditoría. \\
\hline
\textbf{RF-15} &
Extracción de Entidades Nombradas &
El sistema debe identificar entidades nombradas en el texto clasificado como
relevante: (a)~Entidades geográficas (estado, municipio) mediante patrones
georreferenciados; (b)~Víctimas adultas (nombre propio de la madre feminicidada);
(c)~Indicadores de resguardo institucional (DIF, Casa Hogar, tutor). Los nombres
de menores nunca deben ser extraídos ni almacenados (RNF-04). \\
\hline
\end{tabular}
\end{table}

%----------------------------------------------------------------------------------------
\subsection{RF — Subsistema de Agrupamiento Semántico (BERTopic)}
\label{subsec:rf_clustering}
%----------------------------------------------------------------------------------------

\begin{table}[H]
\centering
\caption{Requerimientos Funcionales: Subsistema de Agrupamiento BERTopic}
\label{tab:rf_clustering}
\renewcommand{\arraystretch}{1.4}
\begin{tabular}{|p{1.5cm}|p{3.5cm}|p{8.5cm}|}
\hline
\rowcolor{black}
\textcolor{white}{\textbf{ID}} &
\textcolor{white}{\textbf{Requerimiento}} &
\textcolor{white}{\textbf{Descripción Técnica}} \\
\hline
\textbf{RF-16} &
Reducción de Dimensionalidad UMAP &
El sistema debe reducir los embeddings BETO de 768 a 5 dimensiones mediante UMAP
con parámetros: \texttt{n\_neighbors=15}, \texttt{min\_dist=0.0},
\texttt{metric='cosine'}, \texttt{random\_state=42}. La proyección debe ejecutarse
sobre el corpus completo del ciclo antes de la clasificación individual. \\
\hline
\textbf{RF-17} &
Clustering HDBSCAN &
El sistema debe ejecutar HDBSCAN sobre el espacio UMAP con parámetros:
\texttt{min\_cluster\_size=8}, \texttt{min\_samples=5},
\texttt{cluster\_selection\_method='eom'}, \texttt{prediction\_data=True}.
Los documentos con \texttt{topic\_id=-1} (outliers) deben ser procesados por
la reducción multi-etapa (probabilidades → distribuciones → embeddings) con
objetivo de outliers $\leq 55\%$. \\
\hline
\textbf{RF-18} &
Representación c-TF-IDF de Tópicos &
Para cada clúster identificado, el sistema debe calcular el c-TF-IDF con
\texttt{ngram\_range=(1,3)} y stop words en español, generando los 10 términos
más representativos. La representación final debe refinarse con \texttt{KeyBERTInspired}
(similitud coseno embedding del n-grama vs.\ centroide del clúster). \\
\hline
\textbf{RF-19} &
Contexto de Outlier para Clasificación &
Las noticias con \texttt{topic\_id=-1} tras la reducción de outliers deben recibir
una penalización contextual en la lógica neuro-simbólica: el umbral de BETO se
eleva de 0.50 a 0.65 para estas noticias, requiriendo mayor certeza semántica
para clasificarlas como relevantes. \\
\hline
\end{tabular}
\end{table}

%----------------------------------------------------------------------------------------
\subsection{RF — Subsistema de Persistencia Relacional}
\label{subsec:rf_persistencia}
%----------------------------------------------------------------------------------------

\begin{table}[H]
\centering
\caption{Requerimientos Funcionales: Subsistema de Persistencia}
\label{tab:rf_persistencia}
\renewcommand{\arraystretch}{1.4}
\begin{tabular}{|p{1.5cm}|p{3.5cm}|p{8.5cm}|}
\hline
\rowcolor{black}
\textcolor{white}{\textbf{ID}} &
\textcolor{white}{\textbf{Requerimiento}} &
\textcolor{white}{\textbf{Descripción Técnica}} \\
\hline
\textbf{RF-20} &
Persistencia Transaccional ACID &
Toda inserción de resultados del pipeline (noticia + detección + cluster) debe
ejecutarse dentro de una transacción SQLAlchemy. En caso de excepción en cualquier
operación, se debe ejecutar \texttt{session.rollback()} automático para prevenir
registros huérfanos o parcialmente procesados. \\
\hline
\textbf{RF-21} &
Índice GIN para Full-Text Search &
La columna \texttt{noticias.search\_vector} (\texttt{tsvector}) debe mantenerse
actualizada mediante un trigger SQL que se ejecuta en cada \texttt{INSERT} y
\texttt{UPDATE}, concatenando título (peso A) y contenido (peso B) con el
diccionario Snowball español. El índice GIN debe garantizar consultas FTS en
$O(\log n)$. \\
\hline
\textbf{RF-22} &
API REST de Consulta &
El sistema debe exponer los endpoints \texttt{GET /api/v1/noticias},
\texttt{GET /api/v1/clusters}, \texttt{GET /api/v1/detecciones} con soporte de
paginación (\texttt{page}/\texttt{per\_page}), filtros combinados y respuesta
en \texttt{application/json}. El tiempo de respuesta máximo es de 500~ms para
consultas sobre el índice GIN. \\
\hline
\end{tabular}
\end{table}

%========================================================================================
\section{Requerimientos No Funcionales}
\label{sec:rnf}
%========================================================================================

\begin{table}[H]
\centering
\caption{Requerimientos No Funcionales: Métricas de Calidad y Restricciones}
\label{tab:rnf}
\renewcommand{\arraystretch}{1.4}
\begin{tabular}{|p{1.5cm}|p{3.2cm}|p{8.8cm}|}
\hline
\rowcolor{black}
\textcolor{white}{\textbf{ID}} &
\textcolor{white}{\textbf{Atributo}} &
\textcolor{white}{\textbf{Descripción Técnica y Métrica}} \\
\hline
\textbf{RNF-01} &
Rendimiento: Procesamiento por Lote &
El pipeline completo (recolección + deduplicación + BETO + BERTopic + persistencia)
debe procesar un lote de 600--1{,}000~noticias en un tiempo máximo de 45~minutos
sobre CPU de propósito general (Intel Core i5, 8~GB RAM, sin GPU), garantizando
que el ciclo de 6 horas sea holgadamente completable. La inferencia BETO sobre
lotes de 100~noticias no debe exceder 3~minutos. \\
\hline
\textbf{RNF-02} &
Rendimiento: Consultas en Tiempo Real &
Las consultas FTS desde el dashboard (filtros por eje semántico, rango de fechas,
prioridad) deben retornar resultados en \textless~200~ms sobre una tabla con
$\leq 100{,}000$ registros, garantizado mediante el índice GIN. Las respuestas de
la API REST deben cumplir el SLA de 500~ms al percentil 95. \\
\hline
\textbf{RNF-03} &
Integridad Transaccional (ACID) &
La base de datos PostgreSQL debe operar bajo el modelo ACID con MVCC.
Ante fallo del contenedor \texttt{nna-analyzer} durante un batch, el Write-Ahead
Log debe garantizar que los registros confirmados (\texttt{COMMIT}) persistan y
los no confirmados no corrompan el esquema. El \textit{Rollback} automático debe
ejecutarse en \textless~1~s ante excepción Python. \\
\hline
\textbf{RNF-04} &
Privacidad y Ética (LGDNNA) &
El sistema tiene prohibición absoluta de almacenar, indexar, exportar o renderizar
los nombres propios de los menores afectados. El módulo de extracción de entidades
debe aplicar una lista de exclusión de roles menores (\textit{``hijo(s)'',
``hija(s)'', ``menor(es)''}) que impide la asociación nombre--menor en cualquier
capa del pipeline. Esta restricción opera como constraint a nivel de código, no
como configuración. \\
\hline
\textbf{RNF-05} &
Seguridad: Autenticación (OWASP A07) &
Las contraseñas de acceso al dashboard deben almacenarse exclusivamente como hashes
\textbf{Argon2id} (parámetros: memory=65536~KB, iterations=3, parallelism=4),
resistente a ataques de GPU y side-channel. Las sesiones Flask deben gestionarse
mediante tokens firmados con HMAC-SHA256 (secret\_key de 256~bits, rotación
trimestral). Tras 5 intentos fallidos se bloquea la cuenta por 15~minutos
(OWASP A07:2021 --- Identification and Authentication Failures). \\
\hline
\textbf{RNF-06} &
Seguridad: Protección de API (OWASP A01) &
Todos los endpoints \texttt{/api/v1/*} deben requerir autenticación mediante
token de sesión válido. Las respuestas deben incluir las cabeceras de seguridad:
\texttt{X-Content-Type-Options: nosniff}, \texttt{X-Frame-Options: DENY},
\texttt{Content-Security-Policy}. Las entradas de filtro de consulta deben
ser parametrizadas vía SQLAlchemy ORM (prevención de SQL Injection, OWASP A03). \\
\hline
\textbf{RNF-07} &
Resiliencia: Exponential Backoff &
Ante errores HTTP 429 (Rate Limit) o 503 (Service Unavailable) durante la
recolección, el sistema debe reintentar con espera exponencial: $t_k = t_0 \cdot 2^k$
con $t_0=5$~s y máximo 5 reintentos ($t_{max}=160$~s). Si el quinto reintento
falla, se registra el error en el log y se continúa con la siguiente fuente sin
abortar el ciclo completo. \\
\hline
\textbf{RNF-08} &
Escalabilidad y Portabilidad &
El sistema debe poder desplegarse en cualquier host con Docker Engine $\geq$20.10
mediante \texttt{docker compose up -d} sin modificaciones de código. El esquema
de PostgreSQL debe aplicarse automáticamente mediante \texttt{SQLAlchemy
create\_all()} en el primer arranque. El volumen de procesamiento debe escalar
linealmente al aumentar los workers de Gunicorn o al migrar a un servidor con más
núcleos de CPU. \\
\hline
\textbf{RNF-09} &
Reproducibilidad Científica &
Los modelos empleados (BETO \texttt{dccuchile/bert-base-spanish-wwm-cased},
\texttt{paraphrase-multilingual-MiniLM-L12-v2}) deben fijarse por revisión exacta
en el archivo \texttt{requirements.txt} y en la caché \texttt{models\_cache}
del volumen Docker. Los parámetros de UMAP, HDBSCAN y temperature scaling deben
estar definidos en constantes auditables en \texttt{bertopic\_clustering.py} y
\texttt{semantic\_detector.py}. \\
\hline
\end{tabular}
\end{table}
